% ==============================================================================
% Section 8: Discussion
% ==============================================================================

\section{Discussion}
\label{sec:discussion}

Discussion will interpret the results of the evaluation study in the context of the research questions and the existing literature. The following subsections outline the planned structure.

\subsection{Key Findings}
\label{subsec:key_findings}

This subsection will summarise the principal findings of the study, addressing each research question in turn. RQ1 findings will be contextualised against the reported effect sizes for AR in education \cite{garzonybarragan2020, aldokhny2022}. RQ2 findings will be interpreted through the lens of the ARCS motivational model and compared with engagement metrics reported in gamification studies \cite{deterding2011}. RQ3 findings will be discussed in relation to the ITS effectiveness literature \cite{kulik2016, vanlehn2011} and the role of BKT in adaptive content delivery \cite{corbett1994}.

\subsection{Comparison with Existing Literature}
\label{subsec:comparison_literature}

This subsection will position AR-DeC's results within the broader landscape of AR educational applications, intelligent tutoring systems, and gamified learning platforms. Specific comparisons will be drawn with the findings of Wu et al. \cite{wu2013} on AR in education, Kulik and Fletcher \cite{kulik2016} on ITS effectiveness, and relevant mobile AR learning studies \cite{ibanezmolina2018, parmaxi2020}. The discussion will highlight how the integration of multiple approaches in AR-DeC yields outcomes that differ from studies examining each approach in isolation.

\subsection{Implications for Research and Practice}
\label{subsec:implications}

This subsection will discuss the theoretical and practical implications of the findings. Theoretical implications will address the viability of combining the ARCS model with Bloom's taxonomy as a dual framework for mobile AR learning design. Practical implications will consider the feasibility of deploying AR-DeC-style applications in real educational settings, the role of on-device processing in ensuring accessibility, and design recommendations for practitioners developing similar tools.

\subsection{Limitations}
\label{subsec:limitations}

This subsection will acknowledge the limitations of the study, including: the quasi-experimental design and potential selection bias; the limited sample size and single-institution setting, which constrain generalisability; the four-week intervention duration, which may be insufficient to observe long-term retention effects; the potential novelty effect of the AR technology, which may inflate short-term engagement metrics; device heterogeneity and its impact on user experience; and the reliance on self-report measures for vocabulary knowledge and motivation.
